\scanpage{186}

$(x,\overline{\Sigma}_*)\longmapsto x$, d'où

(301)\quad\diagram{\begin{tikzcd}[ampersand replacement=\&]D_{\Delta_r}^{*}\arrow[r]\&X\end{tikzcd}}

\margin{les images des composantes connexes de $D_{\Delta_r}^{*}$ sont des strates}

qui est une immersion locale. En fait, cet espace sur $X$ ne dépend que de la \inserted{structure de} stratification locale sur $X$ définie par la stratification globale --- on aurait pu le définir directement en termes de la structure locale, à l'aide de la notion de ``drapeaux en germes de strates''. \deleted{Une autre façon de procéder serait de}

\deleted{On aura} \inserted{Dans} le cas d'une stratification associée à un diviseur à cr. norm. sur un esp. analytique, on retrouve les ouverts déjà notés
\[
D_{\Delta_r}^{*}=\coprod_{\substack{\deleted{\unsure{d_*}}\\\unsure{d_r|d_*}\,\unsure{\subset|\in}\,\mathcal P_{r+1}(\mathbb N)}}D_{\unsure{d_r|d_*}}^{*}
\]
de $D_{\Delta_r}$ des deux sections précédentes.

Mais il faudrait pouvoir définir l'analogue des morphismes $\unsure{\tfrac12}$ \unsure{simpliciaux}

\margin{185}
